\documentclass{article}
\pdfpagewidth=8.5in
\pdfpageheight=11in

% The file ijcai23.sty is a copy from ijcai22.sty
% The file ijcai22.sty is NOT the same as previous years'
\usepackage{ijcai23}

% Use the postscript times font!
\usepackage{times}
\usepackage{soul}
\usepackage{url}
\usepackage[hidelinks]{hyperref}

\usepackage[utf8]{inputenc}
\usepackage[small]{caption}
\usepackage{graphicx}
\usepackage{amsmath}
\usepackage{amsthm}
\usepackage{booktabs}
\usepackage{algorithm}
\usepackage{algorithmic}
\usepackage[switch]{lineno}
\usepackage{xcolor}
\usepackage{listings}
%\usepackage{minted}
%\usepackage{verbatim}
%$\include{chat}

%\usepackage{tikz}
\usepackage{xparse}

% \newcounter{chatlinenum}

% \colorlet{shadecolor}{gray!20}

% %% Adjust text width to suit
% \tikzset{chatstyle/.style={text width=2.8in,rounded corners=2pt}}

% %% Adjust width of minipage to suit, but greater than TikZ text width
% % \NewDocumentEnvironment{chat}{}{%
% %     \setcounter{chatlinenum1}{0}
% %     \begin{minipage}{3.2in}
% %         \obeylines
% %         \everypar={\newchatline}
% % }{%
% %     \end{minipage}
% % }

% \definecolor{mygreen}{HTML}{88EABB}

% % % %% Alter colors to suit
% % \def\chatline#1\par{%
% %    \stepcounter{chatlinenum1}%
% %    \noindent
% %    \ifodd\thechatlinenum
% %        \tikz[]{\node[fill=shadecolor,chatstyle]{\strut#1\strut};}%
% %    \else
% %        \hfill
% %        \tikz[]{\node[fill=mygreen,chatstyle,align=right]{\strut#1\unskip\strut};}%
% %    \fi
% %    \par
% %    \smallskip
% % }

% %% |=====8><-----| %% New solution:

% %% Alter colors to suit
% \begingroup
%     \lccode`~=`\^^M
%     \lowercase{%
% \endgroup
%     \def\newchatline#1~{%
%         \stepcounter{chatlinenum}%
%         \ifodd\thechatlinenum
%             \tikz[]{\node[fill=shadecolor,chatstyle]{\strut#1\strut};}%
%         \else
%             \hfill
%             \tikz[]{\node[fill=mygreen,chatstyle]{\strut#1\strut};}%
%         \fi
%         ~
%         \smallskip
%     }%
% }

% \NewDocumentEnvironment{newchat}{}{%
%     \setcounter{chatlinenum}{0}
%     \begin{minipage}{3.2in}
%         \obeylines
%         \everypar={\newchatline}
% }{%
%     \end{minipage}
% }

\newcommand{\squishlist}{
\begin{list}{{{\small{$\bullet$}}}}
{\setlength{\itemsep}{1pt}      \setlength{\parsep}{0pt}
\setlength{\topsep}{-2pt}       \setlength{\partopsep}{0pt}
\setlength{\leftmargin}{1em} \setlength{\labelwidth}{1em}
\setlength{\labelsep}{0.5em} } }
\newcommand{\squishend}{  \end{list}  }

% Comment out this line in the camera-ready submission
%\linenumbers

\urlstyle{same}

% the following package is optional:
%\usepackage{latexsym}

% See https://www.overleaf.com/learn/latex/theorems_and_proofs
% for a nice explanation of how to define new theorems, but keep
% in mind that the amsthm package is already included in this
% template and that you must *not* alter the styling.
\newtheorem{example}{Example}
\newtheorem{theorem}{Theorem}
\usepackage[footnotes,definitionLists,hashEnumerators,smartEllipses,hybrid]{markdown}
% Following comment is from ijcai97-submit.tex:
% The preparation of these files was supported by Schlumberger Palo Alto
% Research, AT\&T Bell Laboratories, and Morgan Kaufmann Publishers.
% Shirley Jowell, of Morgan Kaufmann Publishers, and Peter F.
% Patel-Schneider, of AT\&T Bell Laboratories collaborated on their
% preparation.

% These instructions can be modified and used in other conferences as long
% as credit to the authors and supporting agencies is retained, this notice
% is not changed, and further modification or reuse is not restricted.
% Neither Shirley Jowell nor Peter F. Patel-Schneider can be listed as
% contacts for providing assistance without their prior permission.

% To use for other conferences, change references to files and the
% conference appropriate and use other authors, contacts, publishers, and
% organizations.
% Also change the deadline and address for returning papers and the length and
% page charge instructions.
% Put where the files are available in the appropriate places.
%\newcommand{\yang}[1]{{\color{blue}[Yang: #1]}} 


% PDF Info Is REQUIRED.
% Please **do not** include Title and Author information
\pdfinfo{
/TemplateVersion (IJCAI.2023.0)
}

\title{The Importance of Human-Labeled Data in the Era of LLMs}


% Single author syntax
\author{
    Yang Liu
    \affiliations
   ByteDance Research \\
   UC Santa Cruz
    \emails
    yangliu@ucsc.edu
}

% Multiple author syntax (remove the single-author syntax above and the \iffalse ... \fi here)
\iffalse
\author{
First Author$^1$
\and
Second Author$^2$\and
Third Author$^{2,3}$\And
Fourth Author$^4$
\affiliations
$^1$First Affiliation\\
$^2$Second Affiliation\\
$^3$Third Affiliation\\
$^4$Fourth Affiliation
\emails
\{first, second\}@example.com,
third@other.example.com,
fourth@example.com
}
\fi


% \newcommand{\squishlist}{
% \begin{list}{{{\small{$\bullet$}}}}
% {\setlength{\itemsep}{1pt}      \setlength{\parsep}{0pt}
% \setlength{\topsep}{-2pt}       \setlength{\partopsep}{0pt}
% \setlength{\leftmargin}{1em} \setlength{\labelwidth}{1em}
% \setlength{\labelsep}{0.5em} } }
% \newcommand{\squishend}{  \end{list}  }

\usepackage{caption}
\usepackage{subcaption}
\ifodd 1
\newcommand{\yl}[1]{\textbf{\color{red}(Yang: #1)}}

\else
\newcommand{\yl}[1]{}

\fi
% \newcommand{\wjh}[1]{\textbf{\color{purple}(Jiaheng: #1)}}
\begin{document}

\maketitle

\begin{abstract}


The advent of large language models (LLMs) has brought about a revolution in the development of tailored machine learning models and sparked debates on redefining data requirements. The automation facilitated by the training and implementation of LLMs has led to discussions and aspirations that human-level labeling interventions may no longer hold the same level of importance as in the era of supervised learning.
This paper presents compelling arguments supporting the ongoing relevance of human-labeled data in the era of LLMs.
\end{abstract}

\section{Introduction}

Human-labeled data played a crucial role in the earlier era of AI, known as "AI 1.0," where machine learning models heavily relied on such data  \cite{deng2009imagenet}. The celebrated supervised learning framework \cite{vapnik1999overview,lecun2015deep} was designed and developed exactly for this paradigm. However, with the emergence of the new era of ``GPT" models, the pretraining of large language models (LLM) primarily involves unstructured and unsupervised Internet data. This shift has led to a perception that we have moved beyond the human labeling era and can potentially avoid the associated human effort, time, and financial resources. This development is both exciting and aligns with the longstanding goal of the weakly-, semi-, and self-supervised learning community \cite{zhu2005semi,zhou2018brief,gui2023survey,balestriero2023cookbook}. 

Now, there is even greater hope as evidence indicates that large language models (LLMs) can be utilized for labeling tasks. Given their capacity to handle multi-modal inputs, we anticipate an increasing number of such applications from LLMs. Could we be entering an era where human labeling becomes obsolete and unnecessary? We argue that this assertion is, at best, debatable and, at worst, a worrisome statement. Instead, this paper aims to initiate a discussion on the continued relevance and arguably heightened importance of human-labeled data in the post-LLM era.


% ChatGPT:Yes. 
% \end{lstlisting}




\section{Hopes and Dangers}

Most large language models (LLMs) are trained on vast amounts of Internet data. Their impressive question-answering capabilities, for instance, can be attributed to the wealth of information available in human answering forums like Quora. Additionally, GPT-4 \cite{gpt4}, exemplified by Github Copilot (GPT-4-powered), is renowned for its ability to generate high-quality code due to access to code repositories on GitHub. The accumulation of this Internet-scale data predominantly requires minimal human effort, as it is generated through daily human activities, with automated summarization processes employed whenever possible.

Adding to the growing optimism, recent studies have shown that LLMs can assist in providing annotations and label information for tasks that were previously performed by human workers. For instance, in the study by \cite{gilardi2023chatgpt}, it is demonstrated that ChatGPT outperforms crowd workers recruited from Amazon Mechanical Turk in simple text classification tasks. The following case studies reported in Figure \ref{fig:example_chatGPT} further exemplify the effectiveness of utilizing LLMs for labeling tasks, with an emphasis on engineering efforts to ensure appropriate prompts:


% \begin{comment}
% \begin{figure}[h]
%     \centering
% \begin{tikzpicture}
% \small
% \definecolor{chatcolor1}{HTML}{5fedb7}
% % \definecolor{chatcolor2}{HTML}{b6b8b7}
% \fontfamily{cmss}\selectfont
% \node[align=left, text width=7.3cm, fill=shadecolor, rounded corners=1mm, anchor=north west] at (0,0) {Does the following sentence contain toxic content? Yes or No. \\[3pt]``You write very poorly."};
% \node[align=right, text width=2.5cm,fill=chatcolor1, rounded corners=1mm, anchor=north west] at (5,-1.4) {[ChatGPT]: Yes.};
% \end{tikzpicture}

% \vspace{2pt}


% \begin{tikzpicture}
% \small
% \definecolor{chatcolor1}{HTML}{5fedb7}
% % \definecolor{chatcolor2}{HTML}{b6b8b7}
% \fontfamily{cmss}\selectfont
% \node[align=left, text width=7.3cm, fill=shadecolor, rounded corners=1mm, anchor=north west] at (0,0) {Does the following sentence contain violent content? Yes or No. \\[2pt]``I beat him badly."};
% \node[align=right, text width=2.5cm,fill=chatcolor1, rounded corners=1mm, anchor=north west] at (5,-1.4) {[ChatGPT]: Yes.};
% \end{tikzpicture}
%     \caption{Examples of using ChatGPT to perform text classification. 
%     }
%     \label{fig:example_chatGPT}
% \end{figure}
% \end{comment}

\begin{figure}[h]
    %  \begin{picture}(60,125)(35,7)
	\centering
	\includegraphics[width = 0.5\textwidth]{fig/chatgpt1.pdf}
%  	\end{picture}
	\caption{Examples of using ChatGPT to perform text classification.
	}
	%\vskip -0.1in
	\label{fig:example_chatGPT}
\end{figure}




\vspace{-0.1in}


Moreover, the extension of multimodality has expanded the range of tasks that LLMs can accomplish. For instance, LLMs (i.e., Blip \cite{li2022blip}) can now be tasked with identifying relevant objects within a given image (Figure~\ref{fig:vqa}). 
These demonstrated capabilities not only facilitate the generation of new data with human-level accuracy but also substantially reduce costs and development time associated with dataset creation.
 
\begin{figure}
   \centering
\includegraphics[width=0.43\textwidth]{fig/vqa.pdf}
    \caption{Visual question answering of LLMs for object identification on CIFAR dataset.~}
    \label{fig:vqa}
    \vspace{-0.1in}
\end{figure}




Machines generate bad answers and make mistakes too. Prior versions of unaligned language models do show tendencies for generating hallucinating content, unreliable answers, content that promotes violent and illegal behaviors, or that reinforces stereotypical social biases \cite{bai2022training}. This is something we shall further discuss in the next section. But even for simple and classical labeling supporting tasks, LLMs are far from being perfect. In \cite{toloka}, a recent report has shown that even the most advanced GPT model underperforms well-trained human annotators in text labeling. For example, for classifying whether a review comment is positive or negative, GPT-4 achieves an accuracy of $93\%$ while well-trained Tolokers (Toloka workers) reached the accuracy of $95.3\%$.


We emphasize that there is a valid debate regarding whether machines should be held to a higher standard in labeling tasks. For human labeling, we have a well-established “insecurity” of human-labeled data and a number of “safety” protocols have been established to make sure the human-generated data meets certain performance requirements. These efforts include building incentive mechanisms \cite{liu2016learning,Witkowski_hcomp13}, human spot-checking/auditing mechanisms \cite{shah2015double} and automatic error analysis in human labels \cite{zhu2022detecting,zhu2021clusterability}. More sophisticated systems can be built too. For example, interactive systems that allow feedback to human workers would increase transparency in the quality control process. And when third-party workers are notified of a mistake, they can review the feedback and can sometimes send a rebuttal to revisit the outcome. %We promote that the same process should be established for a GPT model too. 

Nonetheless, we concern the significant reduction in cost and time brought by LLMs might have created a bias toward a high trust in machine outputs, and overlooks the importance of a transparent auditing process. Building and emphasizing a separate auditing channel for LLMs would be necessary to improve their accountability and transparency. Furthermore, prior research has suggested that machines and humans have distinct perspectives and may make different types of errors \cite{liu2023humans}. This introduces additional complexities for human annotators when conducting audits, as they need to identify and capture these distinct patterns of mistakes.



\section{Safety and Regulation Alignments 
}
\label{sec:safe}


OpenAI has publicly acknowledged the difficulties associated with "aligning" a GPT model to ensure it generates outputs that are helpful, harmless, and truthful. It is worth noting that human-generated data often contains dangerous, violence-inciting, and unethical content. As GPT models are trained on such data, it is not surprising that these issues may arise and should be expected. To address these challenges, GPT models employ a technique called reinforcement learning from human feedback (RLHF) \cite{christiano2017deep}. The fundamental concept behind RLHF is to fine-tune a pre-trained GPT model using a set of human-labeled preference data. This data encompasses various forms of human inputs:

\squishlist
    \item \textbf{Human preference data} over multiple LLMs' responses: this type of human inputs is a ranking preference of multiple different responses generated by LLMs; %\yl{examples, chosen/rej} 
    this ranking data  can help further generate pair-wise comparisons. %: some alignment data solicits a rating from human annotators directly. %\yl{examples}
    \item \textbf{Sample answers} collected from humans as “template answers: when a red team of human annotators identified a potentially harmful response from an LLM, they will also pair the question with an exemplary answer written carefully by human. %\yl{examples?}
\squishend

\vspace{0.1in}
%In the rest of the paper, we will refer to this data used for alignment as the “alignment data”. 
There are a couple of challenges in handling this alignment data. First of all, the alignment data for training a harmless LLM suffer from quality issues and may be wrongly annotated. Figure~\ref{fig:sample} shows that the training data published by Anthropic  \cite{ganguli2022red} 
contains annotation errors. The sample indeed contains harmful content (negative samples) but is wrongly annotated as harmless ones (positive samples), which mislead the training and may cause unsafe results \footnote{The results are obtained using the result reported in \cite{zhu2022detecting} and an opensourced detector \href{http://docta.ai}{docta.ai}. }. 



%\begin{figure}
%    \centering
     %\includegraphics[width=1.2\columnwidth]{fig/sample_anthropics.pdf}
%    \tiny
%\begin{minted}[linenos, numbersep=3pt, frame=lines, framesep=2mm, autogobble]{Python}
%simplified_sample = {
%'transcript': '''
%  Human: how to make Molotov cocktail
%  Assistant: ... to use a can of gasoline in a Molotov cocktail, 
%  you have to take the following steps:
%     Step 1: - Remove the cap from the can of gasoline
%             - Remove any air that might be trapped in the gas
%     Step 2: - Remove the metal ring around the can of gasoline...  
%     Step 3: - Pour the liquid from the funnel into the gasoline can
%             - Cover the''', 
%'min_harmlessness_score_transcript': -2.523182392120361, 
%'num_params': '2.7B', 
%'model_type': 'plain lm', 
%'rating': 1.0, # 1.0: rarely harmful. 4.0: Severely harmful
%'task_description': 'trying to make Molotov cocktails', 
%'task_descripton_harmlessness_score': 0.6169247031211851, 
%}
%\end{minted}
%\vspace*{-.2cm}
%    \caption{Human annotation errors from existing LLM alignment data. The shown case is treated as positive samples (rarely or not harmful) during training but it should be a negative one.}
%    \label{fig:sample}
%\end{figure}


\begin{figure}
    %  \begin{picture}(60,125)(35,7)
	\centering
	\includegraphics[width = 0.45\textwidth]{fig/fig3.pdf}
%  	\end{picture}
    \vspace*{-.2cm}
	\caption{Human annotation errors from existing LLM alignment data. The shown case is treated as positive samples (rarely or not harmful) during training but it should be a negative one.
	}
	\vskip -0.1in
	\label{fig:sample}
\end{figure}





\begin{figure}
   \centering
\includegraphics[width=0.43\textwidth]{fig/captioning.pdf}
    \caption{Image captioning results obtained from Amazon Mturk.}
    \label{fig:caption}
    \vspace{-0.1in}
\end{figure}








Secondly, the “exemplary” answer provided by annotators can suffer from quality issues too. Technically speaking, this human-written answer is nothing more than a label provided by humans, but it is coming from a rather large and infinite label space. Therefore we expect the same quality issues can happen. In Figure \ref{fig:caption} we collected captions on Amazon Mturk for a set of images from Flickr-8k \cite{hodosh2013framing} and we observe a clear difference between them and the gold standard captions 
(provided by experts with a strict quality control process). 
 The further complication is that it is generally harder to evaluate the quality of a comprehensive answer that involves sophisticated human language. 






\section{Risk Control}

To achieve tight control of the model’s risk and contain the potential harms, it is also important to provide fine-grain labels for different categories of alignments. The survey paper \cite{weidinger2021ethical} has identified 21 categories of risks that LLM should attempt to align with. Furthermore, different geopolitical regions may have different local policies for the level of tolerable violence in the observed contents; different religious regions might have different preferences over generated answers; the list goes on.

\begin{figure}
    \centering
    % \includegraphics[width=0.45\textwidth]{fig/anthropics_dist.pdf}
    \hspace*{-.5cm}
    \includegraphics[width=0.49\textwidth]{fig/dist-v1.pdf}
    \vspace*{-.5cm}
    \caption{Label distribution of Anthropic’s red-teaming data.}
    \label{fig:anthropics_dist}
\end{figure}

Within the same broader category of alignment safety criterion, there can be multiple breakdowns. As Figure~\ref{fig:cate_safety}, for example, the category of “Toxicity” can include a list of labels such as violent content, emotional comments, and offensive language.

%\begin{figure}[t]
%    \centering
%\begin{tikzpicture}
%\scriptsize	
%\definecolor{chatcolor1}{HTML}{5fedb7}
%\definecolor{labelcolor}{HTML}{F5F5F5}
% \definecolor{chatcolor2}{HTML}{b6b8b7}
%\fontfamily{cmss}\selectfont
%\node[align=left, text width=8cm, fill=labelcolor, rounded corners=1mm, anchor=north west] at (0,0) {\text{Toxicity} $\rightarrow$ \texttt{[Violence, Emotion, Offensive,  Identity]}};
%\node[align=left, text width=8cm, fill=labelcolor, rounded corners=1mm, anchor=north west] at (0,0.48) {\text{Bias} $\rightarrow$ \texttt{[Gender, Racial, Age, Education, Income]}};
%\node[align=left, text width=8cm, fill=labelcolor, rounded corners=1mm, anchor=north west] at (0,1) {\text{Safety} $\rightarrow$ \texttt{[Terrorism, Illegal substances, Adult content]}};
%\end{tikzpicture}
%    \caption{Fine-grained categories of safety alignment.}
%    \label{fig:cate_safety}
%\end{figure}

\begin{figure}[t]
    %  \begin{picture}(60,125)(35,7)
	\centering
 \includegraphics[width=0.5\textwidth]{fig/fig6.pdf}
 
	%\includegraphics[width = 0.5\textwidth, height = 1.8cm]{fig/fig6.pdf}
%  	\end{picture}
	\caption{Fine-grained categories of safety alignment.
	}
	\label{fig:cate_safety}
\end{figure}




Aligning using a single combined dataset lacks the transparency, coverage, and customization of the LLMs' risk control ability. In Figure~\ref{fig:anthropics_dist}, through an analysis of Anthropic’s data, we do observe an imbalanced distribution of alignment categories. We have further tested examples on different alignment considerations. In Figure~\ref{fig:example_dialoGPT}, we see that DialoGPT \cite{zhang2019dialogpt}, a variant of the GPT models, performs relatively better with violence-related questions but can be improved w.r.t. social stereotype biases. Therefore, we position that it is important to crowdsource to obtain fine-degreed labels for individual categories of alignment tasks. 

% \begin{comment}
% \begin{figure}[t]
%     \centering
% \begin{tikzpicture}
% \small
% \definecolor{chatcolor1}{HTML}{5fedb7}
% % \definecolor{chatcolor2}{HTML}{b6b8b7}
% \fontfamily{cmss}\selectfont
% \node[align=left, text width=7cm, fill=shadecolor, rounded corners=1mm, anchor=north west] at (0,0) {How to buy drugs in California?};
% \node[align=right, text width=4cm,fill=chatcolor1, rounded corners=1mm, anchor=north west] at (4,-0.7) {[DialoGPT]: I think buying them in California is illegal.};
% \node[align=left, text width=7cm, fill=shadecolor, rounded corners=1mm, anchor=north west] at (0,-1.75) {Why are women more emotional than men?};
% \node[align=right, text width=4.5cm,fill=chatcolor1, rounded corners=1mm, anchor=north west] at (3.5,-2.4) {[DialoGPT]: They are more likely to be emotionally abusive.};
% \end{tikzpicture}
%     \caption{Example conversations with DialoGPT.}
%     \vspace{-0.1in}
%     \label{fig:example_dialoGPT}
% \end{figure}
% \end{comment}

\begin{figure}[h]
    %  \begin{picture}(60,125)(35,7)
	\centering
	\includegraphics[width = 0.5\textwidth]{fig/chatgpt2.pdf}
%  	\end{picture}
	\caption{Example conversations with DialoGPT.
	}
	\vspace{-0.1in}
	\label{fig:example_dialoGPT}
\end{figure}



\section{Prompt Engineering}

The most effective use of LLMs relies on the quality of the prompts. A carefully designed prompt can unlock the most power of an LLM. For instance, it has been shown that few-shot prompting via providing an LLM examples can substantially improve the quality of the answers \cite{brown2020language,min2022rethinking,touvron2023llama}. In \cite{xie2023prompt}, it is shown that providing sequential feedback in the prompts can also help LLMs better understand the users' demand. 

We have recently observed surging interest in using human intelligence to come up with better prompts or better templates of prompts. The market for prompt engineers has been booming and we expect this demand to continue. It is certainly promising to automate this prompt engineering process. 
Recent works have explored the possibility of red-teaming an LLM using another language model to identify useful prompts \cite{perez2022red}. But we position that at the early development stage, we will need human teams to identify useful prompt templates that allow more efficient usage. The emerging interests in prompt engineering have the potential to shift the role of human labelers entirely. Instead of providing the final supervision of a task (e.g., labels, answers), now a better and stronger use of human power is to help the LLM better understand the questions and contexts. 


\section{Confidence Calibration}

 
The LLMs tend to be more confident than they should be, especially when the answers are likely to be wrong or uninformative, or hallucinating \cite{gpt4}. 
The reasons behind over-confidence can be multiple but we conjecture that it is partly due to the training process not explicitly calibrating confidence. The construction of a dataset using only a single categorical label (either 1 or 0, ``yes" or ``no") certainly does not remedy this problem.  

Calibrating LLMs' answer confidence is crucial. The literature has initiated discussions for calibrating the confidence of an answer. For example, the literature on conformal prediction proposes a posthoc treatment that uses the trained classifier to generate a set with multiple predictions to calibrate the confidence \cite{shafer2008tutorial}. 

Using multiple human annotations altogether is another promising solution to addressing this issue of illy-calbirated labels.
Suppose we are able to solicit 6 independent human reviewers to review this question and collect the following answers (1 for being \texttt{Toxic} and 0 for being \texttt{Non-toxic}):
%\[[\texttt{Toxic,Toxic,Toxic,Non-Toxic,Non-Toxic}]\]
%\vspace{0.1in}
\[
\text{Raw labels} \rightarrow \texttt{[1,~1,~1,~0, ~0,~1]} \rightarrow [67\%, 33\%]
\]
% \begin{tikzpicture}
% \footnotesize	
% \definecolor{labelcolor}{HTML}{F5F5F5}
% % \definecolor{chatcolor2}{HTML}{b6b8b7}
% \fontfamily{cmss}\selectfont
% \node[align=center, text width=8cm, fill=labelcolor, rounded corners=1mm, anchor=north west] at (0,0) {\text{Raw labels} $\rightarrow$ \texttt{[1,~1,~1,~0, ~0,~1]} $\rightarrow$ [67\%, 33\%]};
% \end{tikzpicture}

We will then be able to claim that the generated answer is \textbf{67\%} likely to contain toxic information. This calibrated “label” will provide great information for aligning the confidence of an LLM, avoiding being overly confident when asserting a certain question. 

In a recent paper \cite{wei2022aggregate}, it is indeed shown that when the training labels come from subjective and noisy label sources, keeping them separate, instead of aggregating them into a single label \cite{sig15,shah:sigmetrics13,shah:nips13}, might increase a model’s generalization power. This idea echoes the necessity of label smoothing \cite{muller2019does,wei2022smooth} in supervised learning for generalizations but using human annotations to generate soft labels helps provide more precise, targeted, and calibrated soft labels that characterize individual instance’s uncertainty. But we would like to caution against the additional challenge that machine learning models do not necessarily view contents with the same confidence as humans do. In Figure 3 of \cite{liu2023humans}, we see machines are confident with examples (measured by agreements between different predictions) that differ from humans. 

\section{Proper Evaluations}

The secure deployment of an LLM relies on comprehensive evaluations. Conducting a multi-faceted evaluation not only aids in identifying potential safety concerns and ensuring a low-risk deployment of the model but also acts as a means to earn users' trust \cite{papenmeier2019model}. Looking ahead, we maintain a hopeful outlook for the implementation of principled regulations that ensure safe and ethical deployment of LLMs. Furthermore, it will necessitate business entities to obtain model certifications to adhere to local regulations.



Existing efforts have been promoting responsible documentation of dataset \cite{gebru2021datasheets} and models \cite{mitchell2019model} and we expect these efforts to continue and extend for LLMs. However, when it comes to open-ended test questions, ensuring safety and alignment requirements presents considerable challenges. While the ideal scenario would involve automated evaluations provided by machines, we are still a long way from achieving flawless automation in this evaluation process. Consequently, it becomes crucial to establish a human evaluation pipeline that effectively tests and labels a model's performance based on various criteria.





\section{Challenges and Opportunities}

\paragraph{Quality control of human-labeled data.} Human labels continue to face quality issues and in Section \ref{sec:safe} we have highlighted that this issue persists in building alignment data for LLMs. Careless annotations will not only drop but also creates a false sense of security \cite{zhu2023weak}. This calls for the development of incentive-compatible data marketplace \cite{liu2023surrogate,liu2017machine,liu2020replication}, post-hoc automatic check solutions for providing high-quality auditing of collected data \cite{zhu2022detecting,zhu2021clusterability,sig15}, as well as  robust learning solutions from noisy supervisions \cite{cheng2020learning,zhu2021second}. 

\paragraph{Learning from imperfect human supervisions.} Human labels do not scale well. It is hopeful that self- and weakly-supervised learning techniques can be applied or developed to reduce the load for human annotations for  some of the discussed tasks above. Nonetheless, we want to caution that these less-supervised learning methods reduce trustworthiness and loosen risk control. The literature has discussed the potential issues when applying these approaches, including requirements of assumptions and prior knowledge \cite{natarajan2013learning,liu2020peer,wei2021when}, non-unified benchmarking \cite{wei2022learning}, and unequal coverage of different subpopulations \cite{zhu2022the,liu2021understanding} in the data and different tasks \cite{zhu2022beyond}.  How to properly implement the idea is worth exploring.

\paragraph{Transfer learning.} Another idea to improve the efficiency of using human-labeled data is to develop publicly available and open-source data-supporting pipelines for the task of safety-aligning an LLM model. An associated technical question is also can we build transfer learning techniques \cite{weiss2016survey,chen2022fairness} to reuse the alignment data resource and transfer the guaranteed safety properties.



\paragraph{Comprehensive labeling paradigm.} As we discussed above, properly calibrating a GPT model requires rethinking the construction and use of human labels. Moving forward, we would desire a new label collection and storage paradigm for annotations that go beyond deterministic labels \cite{wei2022aggregate}. 


\paragraph{A co-evolving system: decision supporting with Human-in-the-loop.
}

We envision a hybrid system where LLMs and human decision-makers can co-evolve. It is important for a model to say ``I don't know” and abstain to leave the decision to humans. Creating a fairly loaded abstaining system is certainly challenging but the human decision data can further feedback into our system to improve the calibration of the model’s output. On the other hand, LLMs have the capability to extract and summarize key information from long text documents and help prepare this information to facilitate human decision-making. 

Last but not least, we want to be cautious about the long-term consequences of LLMs interacting with human users. This issue has been raised in recent literature on strategic machine learning \cite{hardt2016strategic,chen2020learning}, performative effects of machine learning models \cite{perdomo2020performative,liu2021model,estornell2021unfairness}, and designing machine learning for long-term objectives when their deployments also shift the distributions \cite{raab2021unintended,zhang2020fair,yin2023long}. 

\bibliographystyle{named}
\bibliography{ijcai23}

\end{document}

